\subsection{
    Процесс ортогонализации Грама-Шмидта, вывод формулы. Построение ортонормированного базиса.
}

Построить ортонормированный базис можно, отталкиваясь от некоторого исходного базиса, при помощи алгоритма, который называют процессом ортогонализации Грама-Шмидта:

\bigbreak

Пусть $f = (\vec{f}_1, \ldots, \vec{f}_n)$ - некоторый базис в евклидовом $n$-мерном пространстве $\mathcal{E}$. Модифицируя этот базис, мы будем строить новый базис $e = (\vec{e}_1, \ldots, \vec{e}_n)$, который будет ортонормированным. Последовательно вычисляем векторы $\vec{g}_1$ и $\vec{e}_1$, $\vec{g}_2$ и $\vec{e}_2$ и т.д. по формулам:

\begin{align*}
    \vec{g}_1 = \vec{f}_1, \quad \quad \quad \quad \quad \quad \quad \quad \quad \quad \quad \quad \quad \quad \quad
    \quad \quad \quad \quad \quad \thinspace \thinspace \thinspace \thinspace \thinspace \vec{e}_1 = \frac{\vec{g}_1}{\norm{\vec{g}_1}};\\
    \vec{g}_2 = \vec{f}_2 - (\vec{f}_2, \vec{e}_1) \cdot \vec{e}_1, \quad \quad \quad \quad \quad \quad \quad \quad \quad \quad
    \quad \quad \quad \quad \thinspace \thinspace \thinspace \thinspace \thinspace \vec{e}_2 = \frac{\vec{g}_2}{\norm{\vec{g}_2}};\\
    \vec{g}_3 = \vec{f}_3 - (\vec{f}_3, \vec{e}_1) \cdot \vec{e}_1 - (\vec{f}_3, \vec{e}_2) \cdot \vec{e}_2, \quad \quad \quad \quad \quad \quad \quad \quad \thinspace \thinspace \thinspace \thinspace \thinspace \vec{e}_3 = \frac{\vec{g}_3}{\norm{\vec{g}_3}};\\
    \ldots \ldots \ldots \ldots \ldots \ldots \ldots
    \ldots \ldots \ldots \ldots
    \quad \quad \quad \quad \quad \quad
    \quad \quad \quad \quad
    \ldots \ldots \ldots \\
    \vec{g}_n = \vec{f}_n - (\vec{f}_n, \vec{e}_1) \cdot \vec{e}_1 - \ldots - (\vec{f}_n, \vec{e}_{n - 1}) \cdot \vec{e}_{n - 1}, \quad \quad \quad \thinspace \thinspace \thinspace \vec{e}_n = \frac{\vec{g}_n}{\norm{\vec{g}_n}}.
\end{align*}

\begin{proof}~

    Рассмотрим индукцию по количеству векторов $n$.
    \begin{enumerate}
        \item При $n = 1$ утверждение очевидно.
        \item Пусть это утверждение выполнено для количества векторов, равного $n$, докажем его для $n + 1$. 
        
        Т.к. утверждение верно для $n$ векторов, то мы можем считать, что векторы $\vec{g}_1, \ldots, \vec{g}_n$ с указанными свойствами уже построены. Построим вектор $\vec{g}_{n + 1}$ в виде 
        $$\vec{g}_{n + 1} = \vec{f}_{n + 1} + \lambda_1 \vec{g}_1 + \ldots + \lambda_n\vec{g}_n.$$ 
        Линейная оболочка векторов $\vec{g}_1, \ldots,  \vec{g}_{n + 1}$ совпадает с $\vec{f}_1, \ldots, \vec{f}_{n + 1}$ при любых $\lambda_i$, поэтому мы будем подбирать коэффициенты $\lambda_i$ так, чтобы выполнялось условие $(\vec{g}_{n + 1},  \vec{g}_i) = 0$ для всех $i = 1, \ldots, n$. Рассмотрим скалярное произведение $$0 = (\vec{g}_{n + 1}, \vec{g}_i) = (\vec{f}_{n + 1}, \vec{g}_i)+ \lambda_1(\vec{g}_1, \vec{g}_i) + \ldots + \lambda_n(\vec{g}_n, \vec{g}_i).$$
        Поскольку $(\vec{g}_j, \vec{g}_i) = 0$ при $j \ne i$ по предположению индукции, то 
        $$0 = (\vec{f}_{n + 1}, \vec{g}_i) + \lambda_i(\vec{g}_i, \vec{g}_i),$$ следовательно 
        $$\lambda_i = -\frac{(\vec{f}_{n + 1}, \vec{g}_i)}{(\vec{g}_i, \vec{g}_i)} \text{ (знаменатель отличен от нуля).}$$ 
        
        Таким образом, чтобы получить вектор $\vec{g}_{n + 1}$, надо из вектора $\vec{f}_{n + 1}$ вычесть его ортогональные проекции на векторы $\vec{g}_1, \ldots, \vec{g}_n$:
        $$\vec{g}_{n + 1} = \vec{f}_{n + 1} -\frac{(\vec{f}_{n + 1}, \vec{g}_1)}{(\vec{g}_1, \vec{g}_1)} \cdot \vec{g}_1 - \ldots -\frac{(\vec{f}_{n + 1}, \vec{g}_n)}{(\vec{g}_n, \vec{g}_n)} \cdot \vec{g}_n.$$
    \end{enumerate}
\end{proof}
